\documentclass[11pt]{article}

\usepackage[margin=1in]{geometry}
\usepackage{amsmath,amssymb}
\usepackage{microtype}

\title{Supplementary Mathematical Details for\\DiffCircaPower}
\author{}
\date{}

\begin{document}
\maketitle

\section{Setup and notation}

Sampling times for individual subjects are denoted by $t_i$, $i=1,\dots,n$.
Distinct sampling times are denoted by $t_b$, $b=1,\dots,B$, with replicate counts $m_b$.
Thus the total sample size is
\begin{equation}
n=\sum_{b=1}^B m_b.
\label{eq:budget}
\end{equation}
We write the harmonic basis vector as
\begin{equation}
f(t)=\begin{bmatrix}1\\ \cos(\omega t)\\ \sin(\omega t)\end{bmatrix}.
\label{eq:basis}
\end{equation}

For evenly spaced sampling over one full cycle $\tau$, let
\begin{equation}
t_b=\frac{(b-1)\tau}{B},\qquad b=1,\dots,B,
\label{eq:even_times}
\end{equation}
so that
\begin{equation}
\theta_b:=\omega t_b=\frac{2\pi(b-1)}{B}.
\label{eq:theta_def}
\end{equation}

We will use the elementary trigonometric identities
\begin{equation}
\sin(2\theta)=2\sin\theta\cos\theta,
\qquad
\cos^2\theta=\frac{1+\cos(2\theta)}{2},
\qquad
\sin^2\theta=\frac{1-\cos(2\theta)}{2}.
\label{eq:trig_identities}
\end{equation}

\section{Design matrix and expansion of \texorpdfstring{$X^\top X$}{X'X}}

For one group and one gene, the design matrix is
\[
X=
\begin{pmatrix}
1 & \cos(\omega t_1) & \sin(\omega t_1)\\
1 & \cos(\omega t_2) & \sin(\omega t_2)\\
\vdots & \vdots & \vdots\\
1 & \cos(\omega t_n) & \sin(\omega t_n)
\end{pmatrix}.
\]
Therefore
\[
X^\top=
\begin{pmatrix}
1 & 1 & \cdots & 1\\
\cos(\omega t_1) & \cos(\omega t_2) & \cdots & \cos(\omega t_n)\\
\sin(\omega t_1) & \sin(\omega t_2) & \cdots & \sin(\omega t_n)
\end{pmatrix}.
\]

Multiplying entry by entry,
\begin{align}
(X^\top X)_{11}
&=\sum_{i=1}^n 1,\label{eq:xtx11}\\
(X^\top X)_{12}
&=\sum_{i=1}^n \cos(\omega t_i),\label{eq:xtx12}\\
(X^\top X)_{13}
&=\sum_{i=1}^n \sin(\omega t_i),\label{eq:xtx13}\\
(X^\top X)_{22}
&=\sum_{i=1}^n \cos^2(\omega t_i),\label{eq:xtx22}\\
(X^\top X)_{23}
&=\sum_{i=1}^n \sin(\omega t_i)\cos(\omega t_i),\label{eq:xtx23}\\
(X^\top X)_{33}
&=\sum_{i=1}^n \sin^2(\omega t_i).\label{eq:xtx33}
\end{align}
By symmetry, the lower-triangular entries are the same as the corresponding upper-triangular ones, so
\begin{equation}
X^\top X=
\begin{pmatrix}
\sum_{i=1}^n1 & \sum_{i=1}^n\cos(\omega t_i) & \sum_{i=1}^n\sin(\omega t_i)\\
\sum_{i=1}^n\cos(\omega t_i) & \sum_{i=1}^n\cos^2(\omega t_i) & \sum_{i=1}^n\sin(\omega t_i)\cos(\omega t_i)\\
\sum_{i=1}^n\sin(\omega t_i) & \sum_{i=1}^n\sin(\omega t_i)\cos(\omega t_i) & \sum_{i=1}^n\sin^2(\omega t_i)
\end{pmatrix}.
\label{eq:xtx_full_entrywise}
\end{equation}

If the design is described in terms of distinct times $t_b$ with replicate counts $m_b$, then for any function $h$,
\begin{equation}
\sum_{i=1}^n h(t_i)=\sum_{b=1}^B m_b\,h(t_b).
\label{eq:replicate_rule}
\end{equation}
Applying this to each entry in \eqref{eq:xtx_full_entrywise} gives
\begin{equation}
X^\top X=\sum_{b=1}^B m_b\,f(t_b)f(t_b)^\top.
\label{eq:xtx_weighted_outer}
\end{equation}
This makes the design interpretation explicit: changing the replicate counts $m_b$ reweights existing rows, whereas introducing new distinct times $t_b$ changes the set of harmonic basis vectors represented in the design.

\section{Balanced evenly spaced sampling}

Now assume a balanced evenly spaced design:
\begin{equation}
t_b=\frac{(b-1)\tau}{B},\qquad m_b\equiv m,\qquad n=Bm,\qquad B\ge 3.
\label{eq:balanced_assumption}
\end{equation}
Then \eqref{eq:replicate_rule} simplifies to
\begin{equation}
\sum_{i=1}^n h(t_i)=m\sum_{b=1}^B h(t_b).
\label{eq:replicate_rule_balanced}
\end{equation}

We now evaluate the entries of \eqref{eq:xtx_full_entrywise} one by one.

\subsection*{The linear cosine and sine sums}

First consider
\[
\sum_{b=1}^B \cos(\theta_b)
=
\sum_{k=0}^{B-1}\cos\!\left(\frac{2\pi k}{B}\right).
\]
Using the finite cosine sum formula,
\[
\sum_{k=0}^{N-1}\cos(a+kd)
=
\frac{\sin(Nd/2)}{\sin(d/2)}
\cos\!\left(a+\frac{(N-1)d}{2}\right),
\]
with $N=B$, $a=0$, and $d=2\pi/B$, we obtain
\begin{align*}
\sum_{b=1}^B \cos(\theta_b)
&=
\frac{\sin\!\left(B\cdot \frac{2\pi}{2B}\right)}{\sin(\pi/B)}
\cos\!\left(\frac{(B-1)\pi}{B}\right)\\
&=
\frac{\sin(\pi)}{\sin(\pi/B)}
\cos\!\left(\frac{(B-1)\pi}{B}\right)\\
&=
0.
\end{align*}
Similarly, using the finite sine sum formula,
\[
\sum_{k=0}^{N-1}\sin(a+kd)
=
\frac{\sin(Nd/2)}{\sin(d/2)}
\sin\!\left(a+\frac{(N-1)d}{2}\right),
\]
again with $N=B$, $a=0$, and $d=2\pi/B$, we get
\begin{align*}
\sum_{b=1}^B \sin(\theta_b)
&=
\sum_{k=0}^{B-1}\sin\!\left(\frac{2\pi k}{B}\right)\\
&=
\frac{\sin(\pi)}{\sin(\pi/B)}
\sin\!\left(\frac{(B-1)\pi}{B}\right)\\
&=
0.
\end{align*}

Therefore
\begin{align}
\sum_{i=1}^n \cos(\omega t_i)
&=
m\sum_{b=1}^B \cos(\theta_b)
=
m\cdot 0
=
0,
\label{eq:linear_cos_sum_subject}\\
\sum_{i=1}^n \sin(\omega t_i)
&=
m\sum_{b=1}^B \sin(\theta_b)
=
m\cdot 0
=
0.
\label{eq:linear_sin_sum_subject}
\end{align}

\subsection*{The sine--cosine cross term}

Using \eqref{eq:trig_identities},
\[
\sin(\theta_b)\cos(\theta_b)=\frac{1}{2}\sin(2\theta_b).
\]
Hence
\begin{align*}
\sum_{b=1}^B \sin(\theta_b)\cos(\theta_b)
&=
\frac{1}{2}\sum_{b=1}^B \sin(2\theta_b)\\
&=
\frac{1}{2}\sum_{k=0}^{B-1}\sin\!\left(\frac{4\pi k}{B}\right).
\end{align*}
Applying the finite sine sum formula with $N=B$, $a=0$, and $d=4\pi/B$ gives
\begin{align*}
\sum_{k=0}^{B-1}\sin\!\left(\frac{4\pi k}{B}\right)
&=
\frac{\sin\!\left(B\cdot \frac{4\pi}{2B}\right)}{\sin(2\pi/B)}
\sin\!\left(\frac{(B-1)2\pi}{B}\right)\\
&=
\frac{\sin(2\pi)}{\sin(2\pi/B)}
\sin\!\left(\frac{2\pi(B-1)}{B}\right)\\
&=
0,
\end{align*}
provided $B\ge 3$ so that $\sin(2\pi/B)\neq 0$. Therefore
\[
\sum_{b=1}^B \sin(\theta_b)\cos(\theta_b)=0,
\]
and hence
\begin{equation}
\sum_{i=1}^n \sin(\omega t_i)\cos(\omega t_i)
=
m\sum_{b=1}^B \sin(\theta_b)\cos(\theta_b)
=
0.
\label{eq:sincos_sum_subject}
\end{equation}

\subsection*{The squared cosine and sine terms}

Using \eqref{eq:trig_identities},
\[
\cos^2(\theta_b)=\frac{1+\cos(2\theta_b)}{2}.
\]
Therefore
\begin{align*}
\sum_{b=1}^B \cos^2(\theta_b)
&=
\sum_{b=1}^B \frac{1+\cos(2\theta_b)}{2}\\
&=
\frac{1}{2}\sum_{b=1}^B 1
+
\frac{1}{2}\sum_{b=1}^B \cos(2\theta_b)\\
&=
\frac{B}{2}
+
\frac{1}{2}\sum_{k=0}^{B-1}\cos\!\left(\frac{4\pi k}{B}\right).
\end{align*}
Now use the finite cosine sum formula with $N=B$, $a=0$, and $d=4\pi/B$:
\begin{align*}
\sum_{k=0}^{B-1}\cos\!\left(\frac{4\pi k}{B}\right)
&=
\frac{\sin\!\left(B\cdot \frac{4\pi}{2B}\right)}{\sin(2\pi/B)}
\cos\!\left(\frac{(B-1)2\pi}{B}\right)\\
&=
\frac{\sin(2\pi)}{\sin(2\pi/B)}
\cos\!\left(\frac{2\pi(B-1)}{B}\right)\\
&=
0,
\end{align*}
again for $B\ge 3$. Hence
\[
\sum_{b=1}^B \cos^2(\theta_b)=\frac{B}{2},
\]
and thus
\begin{equation}
\sum_{i=1}^n \cos^2(\omega t_i)
=
m\sum_{b=1}^B \cos^2(\theta_b)
=
m\cdot \frac{B}{2}
=
\frac{n}{2}.
\label{eq:cos2_sum_subject}
\end{equation}

Similarly,
\[
\sin^2(\theta_b)=\frac{1-\cos(2\theta_b)}{2},
\]
so
\begin{align*}
\sum_{b=1}^B \sin^2(\theta_b)
&=
\sum_{b=1}^B \frac{1-\cos(2\theta_b)}{2}\\
&=
\frac{1}{2}\sum_{b=1}^B 1
-
\frac{1}{2}\sum_{b=1}^B \cos(2\theta_b)\\
&=
\frac{B}{2}
-
\frac{1}{2}\cdot 0
=
\frac{B}{2}.
\end{align*}
Therefore
\begin{equation}
\sum_{i=1}^n \sin^2(\omega t_i)
=
m\sum_{b=1}^B \sin^2(\theta_b)
=
m\cdot \frac{B}{2}
=
\frac{n}{2}.
\label{eq:sin2_sum_subject}
\end{equation}

\subsection*{Final form of \texorpdfstring{$X^\top X$}{X'X}}

Substituting \eqref{eq:linear_cos_sum_subject}, \eqref{eq:linear_sin_sum_subject},
\eqref{eq:sincos_sum_subject}, \eqref{eq:cos2_sum_subject}, and \eqref{eq:sin2_sum_subject}
into \eqref{eq:xtx_full_entrywise} gives
\[
X^\top X=
\begin{pmatrix}
n & 0 & 0\\
0 & n/2 & 0\\
0 & 0 & n/2
\end{pmatrix}.
\]

The condition $B\ge 3$ is necessary here. If $B=2$, then the sampling times are
$\{0,\tau/2\}$, so
\[
\sin(0)=0,
\qquad
\sin(\pi)=0.
\]
Thus the sine column of the design matrix is identically zero, and the design is not full rank.

\end{document}